\chapter{Rust}\label{ch:iv:rust}

The code looks harmless: take a reference to the best bid in a vector of price levels, push a new level, then print the
reference. The compiler refuses it with error E0502, and the candidate is asked two questions: what bug in the equivalent C++
the refusal just prevented, and how to write the function so that it compiles without cloning the book. Rust interviews at
trading and crypto firms test whether the candidate thinks in the language's rules (ownership, borrowing, lifetimes, the traits
that mark what may cross threads) rather than fighting them, and whether the candidate knows where the guarantees stop: at
\texttt{unsafe} and at logic errors. The performance side is One Quant Book~13, chapter~9; every snippet here is compiled by the
chapter's tests with the pinned toolchain and its error code or output asserted.

\section{Ownership and borrowing, and the errors that teach them}

Every value has one owner; assignment of a non-\texttt{Copy} value moves it and the old name can no longer be used (E0382). A
value may be borrowed either by any number of shared references or by exactly one mutable reference at a time (E0499, E0502),
and no reference may outlive what it refers to (E0515, E0597). The rules are the static version of the defects C++ finds with
sanitizers (\cref{ch:iv:cpp}): a shared reference into a vector that is then pushed to is exactly the dangling reference of a
reallocated buffer.

\omcode{code/interviews/23-rust/rust/snippets/borrow_conflict.rs}{1}{6}{A shared borrow of a level held across a push:
error E0502.}\label{lst:iv:rust:borrow}

\begin{method}[Fixing a borrow error without cloning]\label{met:iv:rust:borrow}
\begin{enumerate}
\item Shorten the borrow: copy the value you need out of the reference before mutating (a \texttt{Copy} type costs nothing).
\item Split the borrow: \texttt{split\_at\_mut}, destructuring a struct into its fields, or iterating with \texttt{iter\_mut}.
\item Use indices instead of references when the collection must change underneath.
\item Use the entry API for ``look up, then insert or update'' on maps.
\item Restructure ownership (move the data into the function that mutates it) before reaching for \texttt{RefCell}.
\end{enumerate}
\end{method}

\omcode{code/interviews/23-rust/rust/src/lib.rs}{15}{26}{Two mutable references into one slice, legally: split the
borrow.}\label{lst:iv:rust:split}

\section{Lifetimes in signatures and structs}

A lifetime parameter names the region a reference is valid for, and the compiler checks that callers respect it. When a
function takes two references and returns one, the compiler cannot guess which input the output borrows from (E0106), and the
signature must say. A struct that holds a reference carries a lifetime too, which is how a zero-copy parser returns fields that
borrow from the input buffer instead of allocating.

\omcode{code/interviews/23-rust/rust/src/lib.rs}{37}{60}{A zero-copy field parser: the returned slices borrow from the
buffer for its lifetime \texttt{'a}.}\label{lst:iv:rust:fields}

\section{Traits, generics and dynamic dispatch}

A generic function bounded by a trait is monomorphised: one copy per concrete type, calls resolved and inlined at compile
time. A trait object \texttt{\&dyn Trait} dispatches through a vtable: one copy of the code, an indirect call, and a restriction
that the trait be object-safe (no generic methods, no \texttt{Self} by value in signatures). Drop order is part of the language:
local variables are dropped in reverse order of declaration and struct fields in declaration order, the opposite of C++
members (\cref{iq:iv:rust:7}).

\section{Unsafe, Send and Sync: what the guarantees are and where they stop}

Safe Rust guarantees memory safety and the absence of data races (One Quant Book~13, chapter~9). \texttt{Send} marks a type whose
values may be moved to another thread, \texttt{Sync} one whose shared references may be; the compiler derives them, and a type
such as \texttt{Rc} (a non-atomic reference count) is neither, so sending it to a thread is rejected (E0277). \texttt{unsafe}
blocks let the programmer assert what the compiler cannot check (an index in bounds, a pointer valid); a sound abstraction
keeps every such assertion justified by checks the safe interface enforces. The guarantees stop at logic: deadlocks, wrong
memory orderings in lock-free code that is otherwise safe, and integer overflow in release builds (which wraps) remain the
programmer's.

\omcode{code/interviews/23-rust/rust/src/lib.rs}{87}{96}{An unchecked index made sound by one check at the
boundary.}\label{lst:iv:rust:unsafe}

\section{Question bank}

\begin{interviewq}[$\star$ \iqroles{developer} \iqfirm{crypto firm}]\label{iq:iv:rust:1}
Why does \cref{lst:iv:rust:borrow} not compile, what C++ bug does the refusal prevent, and how do you fix it without cloning the
vector?
\end{interviewq}

\omcode{code/interviews/23-rust/rust/snippets/use_after_move.rs}{1}{5}{Archiving the fills.}\label{lst:iv:rust:moved}

\begin{interviewq}[$\star$ \iqroles{developer} \iqfirm{proprietary firm}]\label{iq:iv:rust:2}
Why does \cref{lst:iv:rust:moved} not compile? Give two fixes and say what each costs.
\end{interviewq}

\omcode{code/interviews/23-rust/rust/snippets/shadowing.rs}{1}{9}{Shadowing.}\label{lst:iv:rust:shadow}

\begin{interviewq}[$\star$ \iqroles{developer} \iqfirm{any}]\label{iq:iv:rust:3}
What does \cref{lst:iv:rust:shadow} print? Is \texttt{px} mutable?
\end{interviewq}

\omcode{code/interviews/23-rust/rust/snippets/two_mut.rs}{1}{6}{Two mutable borrows.}\label{lst:iv:rust:twomut}

\begin{interviewq}[$\star$ \iqroles{developer} \iqfirm{market maker}]\label{iq:iv:rust:4}
Why is \cref{lst:iv:rust:twomut} rejected although the two elements differ, and how do you write it legally?
\end{interviewq}

\omcode{code/interviews/23-rust/rust/snippets/missing_lifetime.rs}{1}{7}{Returning one of two
references.}\label{lst:iv:rust:lifetime}

\begin{interviewq}[$\star\star$ \iqroles{developer} \iqfirm{crypto firm}]\label{iq:iv:rust:5}
Fix \cref{lst:iv:rust:lifetime}. What does the lifetime you add promise the caller?
\end{interviewq}

\omcode{code/interviews/23-rust/rust/snippets/return_local.rs}{1}{8}{Returning a reference to a
local.}\label{lst:iv:rust:local}

\begin{interviewq}[$\star\star$ \iqroles{developer} \iqfirm{proprietary firm}]\label{iq:iv:rust:6}
Why is \cref{lst:iv:rust:local} rejected, and what should the function return instead?
\end{interviewq}

\omcode{code/interviews/23-rust/rust/snippets/drop_order.rs}{1}{20}{Who is dropped when.}\label{lst:iv:rust:drop}

\begin{interviewq}[$\star\star$ \iqroles{developer} \iqfirm{market maker}]\label{iq:iv:rust:7}
What does \cref{lst:iv:rust:drop} print? How does the order of the struct's fields compare with C++?
\end{interviewq}

\begin{interviewq}[$\star\star$ \iqroles{developer} \iqfirm{crypto firm}]\label{iq:iv:rust:8}
Write an iterator over the \texttt{|}-separated fields of a message in a byte buffer that allocates nothing. What lifetime does
each field have?
\end{interviewq}

\begin{interviewq}[$\star\star$ \iqroles{developer, mle} \iqfirm{any}]\label{iq:iv:rust:9}
Count total volume per symbol from a slice of (symbol, volume) pairs with one map lookup per trade.
\end{interviewq}

\begin{interviewq}[$\star\star\star$ \iqroles{developer} \iqfirm{proprietary firm}]\label{iq:iv:rust:10}
A strategy interface is called once per market-data message. Compare a generic function bounded by the trait with a
\texttt{\&dyn} trait object: code size, speed, flexibility. What makes a trait unusable as a trait object?
\end{interviewq}

\omcode{code/interviews/23-rust/rust/snippets/rc_not_send.rs}{1}{8}{Sharing a book with a thread.}\label{lst:iv:rust:rc}

\begin{interviewq}[$\star\star\star$ \iqroles{developer} \iqfirm{crypto firm}]\label{iq:iv:rust:11}
Why is \cref{lst:iv:rust:rc} rejected? Fix it, and explain the difference between \texttt{Send} and \texttt{Sync}.
\end{interviewq}

\begin{interviewq}[$\star\star\star$ \iqroles{developer} \iqfirm{market maker}]\label{iq:iv:rust:12}
Justify the \texttt{unsafe} block of \cref{lst:iv:rust:unsafe}. What would make it unsound, and how would you check it
mechanically?
\end{interviewq}

\begin{interviewq}[$\star\star\star$ \iqroles{developer} \iqfirm{market maker}]\label{iq:iv:rust:13}
Several threads increment a shared counter. Write it without a lock. Which memory ordering do you use, and when would a weaker
or stronger one be wrong?
\end{interviewq}

\begin{omsources}
\item The Rust Reference and the Rustonomicon (for the toolchain pinned by the series, 1.97.1), and the rustc error index
(E0106, E0277, E0382, E0499, E0502, E0515).
\item One Quant Book~13, chapters 9 and~11 (Rust for low latency; lock-free programming).
\end{omsources}
